\ifdefined\gkver\else\def\gkver{student}\fi
\documentclass[\gkver]{gaokaozhenti}
\begin{document}

\chapter{2000年全国新课程}
%% paperType: 理综物理部分

\begin{enumerate}
\item
%% number: 17
%% paperName: 2000年全国高考新课程第 17 题 7 分
%% typeId: 1
%% type: 单选题
%% chapter: 运动和力
%% point: 动量守恒定律
%% method: 无
%% score: 7
%% degree: 750
%% duplicateId: 0
%% body:
光子的能量为 $h\nu$，动量的大小为 $\frac{h\nu}{c}$。如果一个静止的放射性元素的原子在发生 $\gamma$ 衰变时只发出一个 $\gamma$ 光子，则衰变后的原子核\xzanswer{C}

\fourchoices[answer=C]
{仍然静止}
{沿着与光子运动方向相同的方向运动}
{沿着与光子运动方向相反的方向运动}
{可能向任何方向运动}

%% answer: C
%% memo:
\memoanswer{
原子核衰变发出一个 $\gamma$ 光子，动量的大小为 $\frac{h\nu}{c}$，由动量守恒可知，原子核获得的动量也为 $\frac{h\nu}{c}$，方向与光子运动的方向相反，故 C 正确。考生在做这个题的时候，可能会错误认为光子没有质量，发出光子后由 $mv$ 得到光子的动量为零，误选 A。
}

\item
%% number: 18
%% paperName: 2000年全国高考新课程第 18 题 7 分
%% typeId: 1
%% type: 单选题
%% chapter: 气体性质
%% point: 液柱移动定性分析
%% method: 无
%% score: 7
%% degree: 450
%% duplicateId: 0
%% body:
杨氏双缝干涉实验中，下列说法正确的是（$n$ 为自然数，$\lambda$ 为光波波长）\xzanswer{D}

①在距双缝的路程相等的点形成暗条纹

②在距双缝的路程差为 $n\lambda$ 的点形成亮条纹

③在距双缝的路程差为 $n\frac{\lambda}{2}$ 的点形成亮条纹

④在距双缝的路程差为 $\left(n+\frac{1}{2}\right)\lambda$ 的点形成暗条纹

\fourchoices[answer=D]
{①②}
{②③}
{③④}
{②④}

%% answer: D
%% memo:
\memoanswer{
光屏上某点到双缝的路程差为半波长的偶数倍时（波长的整数倍），该点出现亮条纹；光屏上某点到双缝的路程差为半波长的奇数倍时，该点出现暗条纹，所以①③错②④正确，故 ABC 错误 D 正确。
}

\item
%% number: 19
%% paperName: 2000年全国高考新课程第 19 题 7 分
%% typeId: 1
%% type: 单选题
%% chapter: 电路
%% point: 电容
%% method: 无
%% score: 7
%% degree: 550
%% duplicateId: 0
%% body:
图示的电路图中，$C_{2}=2C_{1}$，$R_{2}=2R_{1}$。下列说法正确的是\xzanswer{A}
\onepicture[w=3.4cm]{figs/03.png}

①开关处于断开状态，电容 $C_{2}$ 的电量大于 $C_{1}$ 的电量

②开关处于断开状态，电容 $C_{1}$ 的电量大于 $C_{2}$ 的电量

③开关处于接通状态，电容 $C_{2}$ 的电量大于 $C_{1}$ 的电量

④开关处于接通状态，电容 $C_{1}$ 的电量大于 $C_{2}$ 的电量

\fourchoices[answer=A]
{①}
{④}
{①③}
{②④}

%% answer: A
%% memo:
\memoanswer{
开关处于断开状态时，电容 $C_{1}$ 和 $C_{2}$ 的电压都为电源电动势，由 $C=\frac{Q}{U}$、$C_{2}=2C_{1}$ 可知电容 $C_{2}$ 的电量大于 $C_{1}$ 的电量，故①正确②错误；开关处于接通状态时，电容 $C_{1}$ 和 $R_{2}$ 并联，电容 $C_{2}$ 和 $R_{1}$ 并联，两电容两端的电压等于和它对应并联电阻的电压，由图可知两电阻串联且 $R_{2}=2R_{1}$，所以 $R_{2}$ 两端的电压为 $R_{1}$ 的 2 倍，由 $C=\frac{Q}{U}$ 可知，两个电容上的电量相等，故③④错误，故 A 正确。
}

\item
%% number: 20
%% paperName: 2000年全国高考新课程第 20 题 7 分
%% typeId: 1
%% type: 单选题
%% chapter: 气体性质
%% point: 玻意耳定律
%% method: 无
%% score: 7
%% degree: 750
%% duplicateId: 0
%% body:
下列说法正确的是\xzanswer{C}

\fourchoices[answer=C]
{当两个分子间的距离小于平衡距离时，分子间的作用表现为引力}
{物体的内能是物体中所有分子热运动动能之和}
{只经历等温过程的理想气体，如果压强增加一倍，则其体积减少一半}
{如果没有能量损失，则热机能把从单独一个热源吸收的热量全部转化成机械能}

%% answer: C
%% memo:
\memoanswer{
两个分子间在平衡距离时，分子间作用力之和为零，小于平衡距离表现为斥力，大于平衡距离表现为引力，故 A 错误；物体内所有分子做无规则运动的动能和分子势能的总和叫做物体的内能，故 B 错误；由 $\frac{p_{1}}{p_{2}}=\frac{V_{2}}{V_{1}}$ 可知 C 正确；由热力学第二定律知不可能从单一热源取热，并把它全部变为功而不产生其他任何影响，故 D 错误。
}

\item
%% number: 21
%% paperName: 2000年全国高考新课程第 21 题 7 分
%% typeId: 1
%% type: 单选题
%% chapter: 光学
%% point: 透镜
%% method: 无
%% score: 7
%% degree: 650
%% duplicateId: 0
%% body:
如图：凸透镜 $L$ 的焦距为 $f$，在离透镜 $1.5f$ 处垂直放置一平面镜 $M$。现在焦点 $F$ 处有一物体 $P$，则在透镜另一侧\xzanswer{C}
\onepicture[w=5.2cm]{figs/05.png}

\fourchoices[answer=C]
{不成像}
{距透镜 $2f$ 处成等大、正立的实像}
{距透镜 $2f$ 处成等大、倒立的实像}
{距透镜 $f$ 处成等大、正立的虚像}

%% answer: C
%% memo:
\memoanswer{
$P$ 在平面镜中成像，距离透镜 $2f$，物体 $P$ 发出的光经过平面镜 $M$ 的反射，反射光线进入透镜 $L$ 成像，相当于 $P$ 的像从距离透镜 $2f$ 处发出光，在透镜 $L$ 中成像。由透镜成像公式 $\frac{1}{f}=\frac{1}{u}+\frac{1}{v}$ 可得，$P$ 的像通过透镜成的像为等大倒立的实像，故 C 正确 ABC 错误。
}

\item
%% number: 22
%% paperName: 2000年全国高考新课程第 22 题 0 分
%% typeId: 1
%% type: 单选题
%% chapter: 功和能
%% point: 动能定理
%% method: 无
%% score: 0
%% degree: 450
%% duplicateId: 0
%% body:
如图所示，$OD$ 是水平面，$AB$ 是斜面。初速为 $v_{0}$ 的物体从 $D$ 点出发沿 $DBA$ 滑动到顶点 $A$ 时速度刚好为零。如果斜面改为 $AC$，让该物体从 $D$ 点出发沿 $DCA$ 滑动到 $A$ 点且速度刚好为零，则物体具有的初速度（已知物体与路面之间的动摩擦因数处处相同且不为零）\xzanswer{B}
\onepicture[w=6.0cm]{figs/06.png}

\fourchoices[answer=B]
{大于 $v_{0}$}
{等于 $v_{0}$}
{小于 $v_{0}$}
{取决于斜面的倾角}

%% answer: B
%% memo:
\memoanswer{
设斜面倾角为 $\theta$，动摩擦因数为 $\mu$，物体由 $D$ 点出发，达到 $A$ 端速度为零，它在 $D$ 处的动能一部分转化为 $A$ 处的重力势能，另一部分则克服摩擦力做功生热，由动能定理得 $-\mu mgs_{DB}-\mu mgs_{AB}\cos\theta-mgh=0-\frac{1}{2}mv^{2}$，整理得 $-\mu mgs_{OD}-mgh=0-\frac{1}{2}mv^{2}$，由此可见，速度 $v$ 的大小只与 $h$ 和 $s_{OD}$ 有关，与斜面的倾角没有关系，故 B 正确。
}

\item
%% number: 23
%% paperName: 2000年全国高考新课程第 23 题 20 分
%% typeId: 3
%% type: 填空题
%% chapter: 原子和原子核
%% point: 天然放射性
%% method: 无
%% score: 20
%% degree: 550
%% duplicateId: 0
%% body:
1956 年李政道和杨振宁提出在弱相互作用中宇称不守恒，并由吴健雄用 \ce{^{60}_{27}Co} 放射源进行了实验验证，次年，李、杨二人为此获得诺贝尔物理奖，\ce{^{60}_{27}Co} 的衰变方程是

\[ \ce{^{60}_{27}Co -> ^{A}_{Z}Ni + ^{0}_{-1}e + \bar{\nu}_{e}} \]

其中 $\bar{\nu}_{e}$ 是反中微子，它的电荷为零，静止质量可认为是零。

\begin{enumerate}
	\item
	Co 与 Fe 同周期，它应在周期表的第 \tkanswer{四} 周期。\ce{^{60}_{27}Co} 的核外电子数为 \tkanswer{27}。在上述衰变方程中，衰变产物 \ce{^{A}_{Z}Ni} 的质量数 $A$ 是 \tkanswer{60}，核电荷数 $Z$ 是 \tkanswer{28}。
	
	\item
	在衰变前 \ce{^{60}_{27}Co} 核静止，根据云室照片可以看出，衰变产物 Ni 和 \ce{^{0}_{-1}e} 的运动径迹不在一条直线上。如果认为衰变产物只有 Ni 和 \ce{^{0}_{-1}e}，那么衰变过程将违背 \tkanswer{动量} 守恒定律。
	
	\item
	无水 \ce{CoCl2} 为深蓝色，吸水后变为粉红色的水合物，水合物受热后又变成无水 \ce{CoCl2}，故常在实验室中用作吸湿剂和空气湿度指示剂。
	
	\[ {\ce{CoCl2}}_{\text{深蓝色}}\cdot\ce{ <=>[\text{吸水}][\triangle] }{\ce{CoCl2}}_{\text{粉红色}}\cdot x\ce{H2O} \]
	
	现有无水 \ce{CoCl2} $65\Ug$，吸水后变成 \ce{CoCl2.$x$H2O} $119\Ug$，水合物中 $x$ 的值是 \tkanswer{6}。
	
	\item
	\ce{^{60}_{27}Co} 是典型的 $\gamma$ 放射源，可用于作物诱变育种。我国应用该方法培育出了许多农作物新品种，如棉花高产品种“鲁棉 1 号”，年种植面积曾达到 3000 多万亩，在我国自己培育的棉花品种中栽培面积最大。$\gamma$ 射线处理作物后主要引起 \tkanswer{基因突变}，从而产生可遗传的变异。除 $\gamma$ 射线，用于人工诱变的其它射线还有 \tkanswer{X 射线}、\tkanswer{中子（流）} 和 \tkanswer{紫外线}。
\end{enumerate}

%% answer:
\jdanswer{
\begin{enumerate}
	\item
	四；27；60；28
	
	\item
	动量
	
	\item
	6
	
	\item
	基因突变；X 射线，中子（流），紫外线
\end{enumerate}
}

%% memo:
\memoanswer{
（1）由“质量数 $=$ 质子数 $+$ 中子数”以及“质子数 $=$ 核外电子数”确定 \ce{^{60}_{27}Co} 的核外电子数；由质量数和电荷数守恒确定 \ce{^{A}_{Z}Ni} 中 $A=60$ 和 $Z=28$。

（2）衰变方程隐含着四个守恒关系：质量数守恒、电荷数守恒、动量守恒和能量守恒。

（3）由化学方程式进行计算。

（4）诱变育种的原理是基因突变，$\gamma$ 射线能提高基因突变频率，基因突变的诱变因素有：化学因素、物理因素、生物因素。
}

\item
%% number: 29
%% paperName: 2000年全国高考新课程第 29 题 14 分
%% typeId: 6
%% type: 计算题
%% chapter: 气体性质
%% point: 玻意耳定律
%% method: 无
%% score: 14
%% degree: 550
%% duplicateId: 0
%% body:
如图，某压缩式喷雾器储液桶的容量是 $5.7\times10^{-3}\Umc$。往桶内倒入 $4.2\times10^{-3}\Umc$ 的药液后开始打气，打气过程中药液不会向外喷出。如果每次能打进 $2.5\times10^{-4}\Umc$ 的空气，要使喷雾器内空气的压强达到 4 标准大气压应打气几次？这个压强能否使喷雾器内的药液全部喷完？（设大气压强为 1 标准大气压）
\onepicture[w=5.5cm, align=right]{figs/08.png}

%% answer:
\jdanswer{
18；可以全部喷完
}

%% memo:
\memoanswer{
设标准大气压为 $p_{0}$，药桶中空气的体积为 $V$，打气 $N$ 次以后，喷雾器中的空气压强达到 4 标准大气压，打入的气体在 1 标准大气压下的体积为 $2.5N\times10^{-4}\Umc$，由理想气体状态方程得

\[ p_{0}V+p_{0}\times2.5N\times10^{-4}=4p_{0}V \tag{①} \]

其中

\[ V=5.7\times10^{-3}\Umc-4.2\times10^{-3}\Umc=1.5\times10^{-3}\Umc \]

代入数值，解得 $N=18$ ②，

为方便分析，忽略喷液管中药液产生的压强，当桶中空气完全充满药桶以后，如果空气压强仍然大于大气压强，则药液可以全部喷出，由玻意耳定律得

\[ 4p_{0}V=5.7\times10^{-3}p \tag{③} \]

解得 $p=1.053p_{0}>p_{0}$ ④，

所以，药液可以全部喷出。
}

\item
%% number: 30
%% paperName: 2000年全国高考新课程第 30 题 19 分
%% typeId: 3
%% type: 填空题
%% chapter: 电场
%% point: 带电粒子的直线运动
%% method: 无
%% score: 19
%% degree: 350
%% duplicateId: 0
%% body:
如图所示，直角三角形的斜边倾角为 $30^\circ$，底边 $BC$ 长为 $2L$，处在水平位置，斜边 $AC$ 是光滑绝缘的。在底边中点 $O$ 处放置一正电荷 $Q$。一个质量为 $m$、电量为 $q$ 的带负电的质点从斜面顶端 $A$ 沿斜边滑下，滑到斜边上的垂足 $D$ 时速度为 $v$。

\begin{enumerate}
	\item
	在质点的运动中不发生变化的是\xzanswer{C}
	
	①动能
	
	②电势能与重力势能之和
	
	③动能与重力势能之和
	
	④动能、电势能、重力势能三者之和
	\fourchoices[answer=C]
	{①②}
	{②③}
	{④}
	{②}
	
	\item
	质点的运动是\xzanswer{D}
	\fourchoices[answer=D]
	{匀加速运动}
	{匀减速运动}
	{先匀加速后匀减速的运动}
	{加速度随时间变化的运动}
	
	\item
	该质点滑到非常接近斜边底端 $C$ 点时速率 $v_{C}$ 为多少？沿斜面向下的加速度 $a_{C}$ 为多少？
\end{enumerate}
\onepicture[w=7.0cm, align=right]{\includesvg{figs/09}}

%% answer:
\jdanswer{
\begin{enumerate}
	\item
	C
	
	\item
	D
	
	\item
	$v_{C}=\sqrt{v^{2}+\sqrt{3}gL}$，$a_{C}=\frac{1}{2}g-\frac{\sqrt{3}kQq}{2mL^{2}}$
\end{enumerate}
}

%% memo:
\memoanswer{
[本题 19 分。第（1）问 4 分，第 2 问 4 分，第 3 问 10 分。正确说明电场力不做功，给 3 分；①式 2 分，得出②式给 1 分；③式 2 分，得出④式给 2 分]

（1）带电质点运动的过程中，只有重力和电场力做功，质点具有动能、重力势能以及电势能，由功能关系知这三种能量之和始终为定值，故 C 正确。

（2）带电质点运动的过程中，所受的重力和斜面的支持力保持不变，但电场力发生变化，故合力的大小一直发生变化，所以质点做加速度随时间变化的运动，故 D 正确。

（3）因 $BD=\frac{BC}{2}=BO=OC=OD$，则 $B$、$C$、$D$ 三点在以 $O$ 为圆心的同一圆周上，是 $O$ 点处点电荷 $Q$ 产生的电场中的等势点，所以，$q$ 由 $D$ 到 $C$ 的过程中电场力做功为零。由机械能守恒定律：

\[ mgh=\frac{1}{2}mv_{C}^{2}-\frac{1}{2}mv^{2} \tag{①} \]

由几何关系得

\[ h=BD\sin60^\circ=BC\sin30^\circ\sin60^\circ=\frac{\sqrt{3}}{2}L \]

联立解得

\[ v_{C}=\sqrt{v^{2}+\sqrt{3}gL} \tag{②} \]

质点在 $C$ 点受三个力的作用：电场力 $F=\frac{kQq}{L^{2}}$，方向由 $C$ 指向 $O$ 点；重力 $mg$，方向竖直向下；支持力 $N$，方向垂直于斜面向上。

根据牛顿定律，有 $mg\sin\theta-F\cos\theta=ma_{C}$ ③，

\[ mg\sin30^\circ-\frac{kQq}{L^{2}}\cos30^\circ=ma_{C} \]

解得 $a_{C}=\frac{1}{2}g-\frac{\sqrt{3}kQq}{2mL^{2}}$ ④。
}

\item
%% number: 31
%% paperName: 2000年全国高考新课程第 31 题 15 分
%% typeId: 4
%% type: 实验题
%% chapter: 电路
%% point: 电路实验
%% method: 无
%% score: 15
%% degree: 350
%% duplicateId: 0
%% body:
电阻 $R_{1}$，$R_{2}$，$R_{3}$ 连结成图示的电路，放在一个箱中（虚框所示），箱面上有三个接线柱 $A$、$B$、$C$。请用多用表和导线设计一个实验，通过在 $A$、$B$、$C$ 的测量，确定各个电阻的阻值。要求写出实验步骤并用所测值表示电阻 $R_{1}$、$R_{2}$、$R_{3}$。
\onepicture[w=5.5cm, align=right]{figs/10.png}

%% answer:
\jdanswer{
解法 1：$R_{1}=\frac{2xyz}{xy+yz-xz}$，$R_{2}=\frac{2xyz}{yz-xy+xz}$，$R_{3}=\frac{2xyz}{xy+xz-yz}$。

解法 2：

$R_{1}=\frac{1}{2(y+z-x)}\left[2(xy+yz+zx)-x^{2}-y^{2}-z^{2}\right]$，

$R_{2}=\frac{1}{2(x+z-y)}\left[2(xy+yz+zx)-x^{2}-y^{2}-z^{2}\right]$，

$R_{3}=\frac{1}{2(x+y-z)}\left[2(xy+yz+zx)-x^{2}-y^{2}-z^{2}\right]$。
}

%% memo:
\memoanswer{
（本题 15 分。实验步骤占 3 分；能正确写出方程式①、②、③的给 6 分，每错 1 个扣 2 分；正确得出 $R_{1}$、$R_{2}$、$R_{3}$ 的表达式再给 6 分，每错一个扣 2 分）

解法 1：比较灵活地用导线短接电路进行测量。实验步骤如下：

①用导线连结 $BC$，测出 $A$、$B$ 两点间的电阻值 $x$；

②用导线连结 $AB$，测出 $B$、$C$ 两点间的电阻值 $y$；

③用导线连结 $AC$，测出 $B$、$C$ 两点间的电阻值 $z$。

则有 $\frac{1}{x}=\frac{1}{R_{1}}+\frac{1}{R_{2}}$ ①，$\frac{1}{y}=\frac{1}{R_{2}}+\frac{1}{R_{3}}$ ②，$\frac{1}{z}=\frac{1}{R_{1}}+\frac{1}{R_{3}}$ ③，

联立①、②两式得 $\frac{1}{x}-\frac{1}{y}=\frac{1}{R_{1}}-\frac{1}{R_{3}}$ ④，

联立③、④两式得 $R_{1}=\frac{2xyz}{xy+yz-xz}$，

同理可解得 $R_{2}=\frac{2xyz}{yz-xy+xz}$，$R_{3}=\frac{2xyz}{xy+xz-yz}$。

解法 2：不用导线短接，直接测量。实验步骤如下：

①测出 $A$、$B$ 两点间的电阻值 $x$；

②测出 $B$、$C$ 两点间的电阻值 $y$；

③测出 $A$、$C$ 两点间的电阻值 $z$；

则有：$\frac{1}{x}=\frac{1}{R_{1}}+\frac{1}{R_{2}+R_{3}}$ ①，

$\frac{1}{y}=\frac{1}{R_{2}}+\frac{1}{R_{1}+R_{3}}$ ②，

$\frac{1}{z}=\frac{1}{R_{3}}+\frac{1}{R_{1}+R_{2}}$ ③，

解得：

$R_{1}=\frac{1}{2(y+z-x)}\left[2(xy+yz+zx)-x^{2}-y^{2}-z^{2}\right]$，

$R_{2}=\frac{1}{2(x+z-y)}\left[2(xy+yz+zx)-x^{2}-y^{2}-z^{2}\right]$，

$R_{3}=\frac{1}{2(x+y-z)}\left[2(xy+yz+zx)-x^{2}-y^{2}-z^{2}\right]$。
}

\end{enumerate}

\end{document}
